\documentclass[12pt]{exam}

\usepackage[brazil]{babel}   
\usepackage[utf8]{inputenc}  
%\usepackage[T1]{fontenc}
\usepackage{amsfonts,amsmath,amssymb,latexsym} 
\usepackage{float}
\usepackage{graphicx}
\graphicspath{ {images/} }

\def\cQ{\mathcal{Q}}
\def\ve{\varepsilon}

\def\emptyset{\varnothing}



\title{Lista 03 de Linguagens Formais e Autômatos}
\date{}
%\date{2º Período de 2023}
\author{Turma do 3º ano}

\begin{document} 

\maketitle

\begin{description}

\item[Operadores de Expressões Regulares]:
\begin{itemize}
\item A união de duas linguagens $L$ e $M$ denotada por $L\cup M$ 
é o conjunto de strings que estão em $L$ ou em $M$
\item A concatenação de duas linguagens $L$ e $M$ denotada por $LM$
é o conjunto de strings que podem ser formados concatenando-se qualquer string de $L$ 
com qualquer string de $M$, 
nesta ordem
\item A potência $L^i$, para $i=0,1,2,3, \dots$; é a concatenação de $L$ com $L^{i-1}$, se $i>0$;
e $L^0 = \{\ve\}$ para $i=0$.
\item O fechamento de uma linguagem $L$, denotado por $L^*$ é a união
\[\bigcup_{i=0}^\infty L^i\]
\end{itemize}

\item[Definição de uma Expressão Regular (ER)]:

\begin{itemize}
\item $\emptyset$ é uma ER; $L(\emptyset) = \emptyset$.
\item $\ve$ é uma ER; $L(\ve) = \{\ve\}$.
\item \textbf{a} é uma ER, $a\in \Sigma$; $L($\textbf{a}$) = \{a\}$.
\item Se $E$ e $F$ são ER, então $E+F$ é ER; $L(E+F) = L(E)\cup L(F)$.
\item Se $E$ e $F$ são ER, então $EF$ é ER; $L(EF) = L(E)L(F)$.
\item Se $E$ é ER, então $E^*$ é ER; $L(E^*) = L(E)^*$.
\item Se $E$ é ER, então $(E)$ é ER; $L((E)) = L(E)$.
\end{itemize}



\end{description}

\vspace{3em}

\break

\begin{questions}



\question Desenvolva ER que gerem as seguintes linguagens sobre o alfabeto $\Sigma = \{a,b\}$:

\begin{parts}

\part $\{w~|~w$ a primeira letra é $a$ $\}$

\part $\{w~|~w$ a última letra é $b$ $\}$

\part $\{w~|~w$ a primeira e a última letra é $b$ $\}$

\part $\{w~|~w$ a penúltima letra é $b$ $\}$

\part $\{w~|~w$ a antipenúltima letra é $b$ $\}$

\part $\{w~|~w $ tem, pelo menos, um par de $a$ como substring$\}$

\part $\{w~|~w $ tem, exatamente, um par de $a$ como substring, podendo ter mais $a$s sem estar ao lado de outro $a$.$\}$

\part $\{w~|~w $ todo $a$ sempre ocorre em par$\}$

\part $\{w~|~w $ tem uma quantidade par de $a$s$\}$

\part $\{w~|~w $ possui $aba$ como substring$\}$

\end{parts}





\end{questions}

\end{document}

